\documentclass[10pt]{article} 

\usepackage[UTF8]{inputenc}
\usepackage[french]{babel}
\usepackage [T1]{fontenc}
\usepackage{sectsty}
\usepackage{graphicx}
\title{Mon document}
\author {Bhouvan}
\sectionfont{\centering}
\renewcommand{\thesection}{\Roman{section}}
\begin{document}


\section{Module1}
\subsection{TP1}
\begin{enumerate}
\item $0.1f$ n'est pas exact parce que sa convertion en base 2 donne un nombre de chiffre infini apres la virgule.
\item Signe, mantisse et exposant de 1.0f
   \begin{figure}[hbtp]
	\caption{img 1}
	\centering
	\includegraphics[scale=1]{C:/Users/DESKTOP-G4VR92/Downloads/1.0f.PNG}
	\end{figure}
\item identifier $+$Inf 	
   \begin{figure}[hbtp]
	\caption{img 2}
	\centering
	\includegraphics[height=4cm, width=5cm, keepaspectratio=false]{C:/Users/DESKTOP-G4VR92/Downloads/inf.PNG}
	\end{figure}

\item identifier NaN 
	\begin{figure}[hbtp]
	\caption{img 3}
	\centering
	\includegraphics[height=4cm, width=5cm, keepaspectratio=false]{C:/Users/DESKTOP-G4VR92/Downloads/nan.PNG}
	\end{figure}
\newpage
\item Bits   de $+0.0$ et $-0.0$. \\ On peut dire que les deux nombres on la meme mantisse et le meme exposant mais leur bits de signe sont different etant respectivement egaux a $0$ et $1$ pour $+0.0$ et $-0.0$
	\begin{figure}[hbtp]
	\caption{img 3}
	\centering
	\includegraphics[height=4cm, width=5cm, keepaspectratio=false]{C:/Users/DESKTOP-G4VR92/Downloads/bits.PNG}
	\end{figure}
	
	
	 
		
\end{enumerate}






\end{document}